\chapter{核心代码}

本附录对应正文第 6 章的代码规范说明。完整代码位于项目目录中的 \texttt{notebooks/}、\texttt{src/} 和 \texttt{scripts/}，此处只列出核心结构和关键片段。

\section{核心源码文件}

\begin{itemize}
  \item \texttt{src/config.py}：路径、随机种子、目标变量、泄漏字段和聚类输入白名单。
  \item \texttt{src/data\_utils.py}：数据下载、读取、合规性检查和摘要输出。
  \item \texttt{src/feature\_engineering.py}：工程特征生成、任务级特征筛选和预处理管道。
  \item \texttt{src/model\_utils.py}：分类、回归、聚类评估指标与结果保存。
  \item \texttt{src/visualization.py}：EDA、模型评估、聚类画像和 PCA 图绘制。
\end{itemize}

\section{核心代码片段}

\begin{lstlisting}[language=Python,caption={统一随机种子与模型配置},label={lst:appendix-random-state}]
RANDOM_STATE = 42

regressors = {
    "linear_regression": LinearRegression(),
    "ridge": Ridge(),
    "random_forest": RandomForestRegressor(random_state=RANDOM_STATE),
    "gradient_boosting": GradientBoostingRegressor(random_state=RANDOM_STATE),
}
\end{lstlisting}

\begin{lstlisting}[language=Python,caption={聚类输入白名单示意},label={lst:appendix-clustering-whitelist}]
CLUSTERING_FEATURES = [
    "device_hours_per_day", "phone_unlocks", "notifications_per_day",
    "social_media_mins", "study_mins", "physical_activity_days",
    "sleep_hours", "sleep_quality", "social_media_hours",
    "study_hours", "notifications_per_device_hour",
    "unlocks_per_device_hour", "device_to_sleep_ratio",
    "activity_sleep_interaction", "social_to_study_ratio",
]
\end{lstlisting}
